\section{Conclusion}
In this report, the advection equation was studied through the Leapfrog scheme and the Upwind scheme. These schemes were studied for different Courant numbers and compared with the analytical solution. A Courant number of $\mu=1.1$ proved to give unstable results from both schemes, while a Courant number of $\mu=1$ gave almost identical results to the analytical solution. However, due to the nature of more complex equations, a Courant number of exactly $\mu=1$ is often unrealistic. Therefore, the Courant number of $\mu=0.9$ is more interesting to study. For a cosine wave, $\mu=1$ proved to give very good results for both schemes, although the Upwind scheme showcased diffusive properties over time. Both schemes showcase a slower phase speed than the theoretical one, although this was a very slight error. For a cosine pulse, the Courant number of $\mu=0.9$ gave diffusive properties for the Upwind scheme and a visible computational mode for the Leapfrog scheme. However, applying a Robert-Asselin parameter to the Leapfrog scheme proved to remove this computational mode.

Overall, this lab has showcased the powerful tools of numerical models through the Leapfrog and Upwind schemes. The advection equation is an equation that can be solved analytically, but the results in this report showcase that these schemes can be used to study even more difficult equations and give reasonable results.